\documentclass{article}
\usepackage{geometry}
\geometry{margin=1in}
\usepackage{array}
\usepackage{enumitem}
\usepackage{booktabs}

\begin{document}

\title{Earth Science: Minerals and Rocks Study Guide}
\author{6th Grade Science}
\date{}
\maketitle

% \tableofcontents
% \newpage

\section{Fundamental Concepts}

\subsection{What is a Mineral?}
A \textbf{mineral} is a naturally occurring, inorganic solid that has:
\begin{itemize}
    \item A crystal structure
    \item A definite chemical composition
    \item Formed through geological processes
\end{itemize}

\subsection{What is a Rock?}
A \textbf{rock} is a naturally formed solid made of one or more minerals that make up the Earth's crust.

\subsection{Key Differences}
\begin{tabular}{@{}p{0.45\textwidth} p{0.45\textwidth}@{}}
\textbf{Minerals} & \textbf{Rocks} \\
\midrule
Single, pure substance & Mixture of minerals \\
Definite chemical formula & Variable composition \\
Crystalline structure & May or may not show crystals \\
Examples: quartz, calcite, gold & Examples: granite, sandstone, marble \\
\end{tabular}

\section{Mineral Identification Properties}

\subsection{Physical Properties for Identification}

\begin{enumerate}
    \item \textbf{Color}
    \begin{itemize}
        \item Quick visual clue but often unreliable alone
        \item Same mineral can have different colors
        \item Different minerals can have the same color
    \end{itemize}
    
    \item \textbf{Streak}
    \begin{itemize}
        \item Color of mineral's powder when rubbed on streak plate
        \item More reliable than surface color
        \item Example: Pyrite has golden color but black streak
    \end{itemize}
    
    \item \textbf{Luster}
    \begin{itemize}
        \item How mineral reflects light
        \item Types: metallic, glassy, pearly, dull, waxy
        \item Example: Gold has metallic luster, quartz has glassy luster
    \end{itemize}
    
    \item \textbf{Hardness}
    \begin{itemize}
        \item Measure of how easily mineral can be scratched
        \item Tested using Mohs Hardness Scale (1-10)
        \item 1 = softest (talc), 10 = hardest (diamond)
    \end{itemize}
    
    \item \textbf{Cleavage vs. Fracture}
    \begin{itemize}
        \item \textbf{Cleavage:} Breaks along flat, planar surfaces
        \item \textbf{Fracture:} Breaks with irregular, uneven surfaces
        \item Example: Mica has perfect cleavage, quartz has conchoidal fracture
    \end{itemize}
    
    \item \textbf{Density}
    \begin{itemize}
        \item Mass per unit volume
        \item "Heft test" - how heavy it feels for its size
        \item Example: Gold feels much heavier than pyrite of same size
    \end{itemize}
\end{enumerate}

\subsection{Mohs Hardness Scale}
\begin{tabular}{@{}p{0.15\textwidth} p{0.25\textwidth} p{0.5\textwidth}@{}}
\textbf{Scale} & \textbf{Mineral (Symbol)} & \textbf{Test} \\
\midrule
1 & Talc & Can be scratched by fingernail \\
2 & Gypsum & Can be scratched by fingernail \\
3 & Calcite (CaCO$_3$) & Can be scratched by copper penny \\
4 & Fluorite (CaF$_2$) & Can be scratched by iron nail \\
5 & Apatite & Can be scratched by glass \\
6 & Feldspar & Can scratch glass \\
7 & Quartz (SiO$_2$) & Can scratch glass easily \\
8 & Topaz & Can scratch quartz \\
9 & Corundum (Al$_2$O$_3$) & Can scratch topaz \\
10 & Diamond (C) & Can scratch everything \\
\end{tabular}

\section{Special Tests for Mineral Identification}

\subsection{Additional Identification Methods}
\begin{itemize}
    \item \textbf{Acid Test:} Calcite fizzes when hydrochloric acid is applied
    \item \textbf{Magnetism:} Magnetite is attracted to magnets
    \item \textbf{Taste:} Halite (rock salt) tastes salty
    \item \textbf{Smell:} Sulfur has a distinctive smell
\end{itemize}

\subsection{Common Mineral Examples}
\begin{tabular}{@{}p{0.2\textwidth} p{0.3\textwidth} p{0.4\textwidth}@{}}
\textbf{Mineral} & \textbf{Key Properties} & \textbf{Uses} \\
\midrule
Quartz (SiO$_2$) & Glassy luster, hardness 7, white streak & Glass, electronics \\
Calcite (CaCO$_3$) & Reacts with acid, hardness 3 & Construction, antacids \\
Gold (Au) & Metallic luster, soft, golden streak & Jewelry, electronics \\
Pyrite (FeS$_2$) & Metallic luster, hard, black streak & Fool's gold, sulfur source \\
Mica & Perfect cleavage, flexible sheets & Electronics, insulation \\
\end{tabular}

\section{Weathering and Erosion}

\subsection{Key Processes}
\begin{enumerate}
    \item \textbf{Weathering}
    \begin{itemize}
        \item Chemical and physical breakdown of rock at Earth's surface
        \item Creates sediment (small pieces of rock)
    \end{itemize}
    
    \item \textbf{Erosion}
    \begin{itemize}
        \item Movement of weathered material by water, wind, or ice
        \item Transports sediment to new locations
    \end{itemize}
    
    \item \textbf{Deposition}
    \begin{itemize}
        \item Laying down of sediment in new locations
        \item Forms layers that may become sedimentary rock
    \end{itemize}
\end{enumerate}

\section{The Three Types of Rocks}

\subsection{Igenous Rocks}
\begin{itemize}
    \item \textbf{Formation:} From cooling of molten rock (magma or lava)
    \item \textbf{Location:} Magma cools underground, lava cools on surface
    \item \textbf{Examples:} Granite, basalt, obsidian
    \item \textbf{Characteristics:} May have crystals, no fossils
\end{itemize}

\subsection{Sedimentary Rocks}
\begin{itemize}
    \item \textbf{Formation:} From compaction and cementation of sediments
    \item \textbf{Process:} Weathering → Erosion → Deposition → Compaction → Cementation
    \item \textbf{Examples:} Sandstone, limestone, shale
    \item \textbf{Characteristics:} May contain fossils, layered structure
\end{itemize}

\subsection{Metamorphic Rocks}
\begin{itemize}
    \item \textbf{Formation:} From heat and pressure on existing rocks
    \item \textbf{Process:} Changes mineral structure without melting
    \item \textbf{Examples:} Marble, slate, gneiss
    \item \textbf{Characteristics:} May have foliation (banding), no fossils
\end{itemize}

\section{The Rock Cycle}
\subsection{How Rocks Change}
The rock cycle shows how rocks are continuously formed, changed, and reformed:
\begin{enumerate}
    \item \textbf{Igneous to Sedimentary:} Weathering and erosion break down igneous rocks into sediments, which are then compacted and cemented to form sedimentary rocks.
    \item \textbf{Sedimentary to Metamorphic:} Heat and pressure deep within the Earth change sedimentary rocks into metamorphic rocks.
    \item \textbf{Metamorphic to Igneous:} Metamorphic rocks can melt into magma, which cools and solidifies to form igneous rocks.
    \item \textbf{Any rock type can become any other rock type} through various processes in the cycle.
\end{enumerate}

\section{Fossils}

\subsection{What are Fossils?}
\begin{itemize}
    \item Preserved traces, imprints, or remains of ancient organisms
    \item Most commonly found in sedimentary rocks
    \item Provide evidence of past life and environments
\end{itemize}

\subsection{Types of Fossils}
\begin{itemize}
    \item \textbf{Body Fossils:} Actual remains (bones, shells)
    \item \textbf{Trace Fossils:} Evidence of activity (footprints, burrows)
    \item \textbf{Imprint Fossils:} Impressions in rock
\end{itemize}

\section{Common Misconceptions and Pitfalls}

\subsection{What is NOT a Mineral?}
\begin{itemize}
    \item \textbf{Living materials:} Ivory, amber, pearls
    \item \textbf{Man-made materials:} Plastic, glass, cement
    \item \textbf{Volcanic glass (obsidian):} Lacks crystalline structure
    \item \textbf{Rocks:} Mixtures of minerals (like sandstone)
\end{itemize}

\subsection{Identification Mistakes to Avoid}
\begin{itemize}
    \item Relying only on color for identification
    \item Confusing similar-looking minerals (gold vs. pyrite)
    \item Not testing multiple properties
    \item Ignoring streak color vs. surface color
\end{itemize}

\section{Study Tips and Assessment}

\subsection{Key Points to Remember}
\begin{enumerate}
    \item Always test multiple properties for mineral identification
    \item Streak is more reliable than color
    \item Hardness tests help distinguish similar minerals
    \item Rocks are made of minerals, minerals are pure substances
    \item The rock cycle is continuous and ongoing
\end{enumerate}
 
\subsection{Lab Activities}
\begin{itemize}
    \item Mineral identification stations with unknown samples
    \item Streak tests using streak plates
    \item Hardness tests using common objects
    \item Acid tests on calcite samples
    \item Rock cycle modeling with clay or playdough
\end{itemize}

\subsection{Assessment Focus}
\begin{itemize}
    \item Understanding mineral definition (5 key characteristics)
    \item Proper use of identification tests
    \item Distinguishing between minerals and rocks
    \item Explaining the rock cycle
    \item Understanding weathering, erosion, and deposition
\end{itemize}

\end{document}
